\documentclass[a4paper,11pt,reqno]{amsart}

\usepackage[utf8]{inputenc}
\usepackage[foot]{amsaddr}
\usepackage{amsmath,amsfonts,amssymb,amsthm,mathrsfs,bm}
\usepackage[margin=0.95in]{geometry}
\usepackage{color}
\usepackage[dvipsnames]{xcolor}

\input{toc-config.tex}

\usepackage{mathtools,enumerate,mathrsfs,graphicx}
\usepackage{epstopdf}
\usepackage{hyperref}

\usepackage{latexsym}


\definecolor{CommentGreen}{rgb}{0.0,0.4,0.0}
\definecolor{Background}{rgb}{0.9,1.0,0.85}
\definecolor{lrow}{rgb}{0.914,0.918,0.922}
\definecolor{drow}{rgb}{0.725,0.745,0.769}

\usepackage{listings}
\usepackage{textcomp}
\lstloadlanguages{Matlab}%
\lstset{
    language=Matlab,
    upquote=true, frame=single,
    basicstyle=\small\ttfamily,
    backgroundcolor=\color{Background},
    keywordstyle=[1]\color{blue}\bfseries,
    keywordstyle=[2]\color{purple},
    keywordstyle=[3]\color{black}\bfseries,
    identifierstyle=,
    commentstyle=\usefont{T1}{pcr}{m}{sl}\color{CommentGreen}\small,
    stringstyle=\color{purple},
    showstringspaces=false, tabsize=5,
    morekeywords={properties,methods,classdef},
    morekeywords=[2]{handle},
    morecomment=[l][\color{blue}]{...},
    numbers=none, firstnumber=1,
    numberstyle=\tiny\color{blue},
    stepnumber=1, xleftmargin=10pt, xrightmargin=10pt
}

\numberwithin{equation}{section}
\synctex=1

\hypersetup{
    unicode=false, pdftoolbar=true, 
    pdfmenubar=true, pdffitwindow=false, pdfstartview={FitH}, 
    pdftitle={ELE2024 Coursework}, pdfauthor={A. Author},
    pdfsubject={ELE2024 coursework}, pdfcreator={A. Author},
    pdfproducer={ELE2024}, pdfnewwindow=true,
    colorlinks=true, linkcolor=red,
    citecolor=blue, filecolor=magenta, urlcolor=cyan
}


% CUSTOM COMMANDS
\renewcommand{\Re}{\mathbf{re}}
\renewcommand{\Im}{\mathbf{im}}
\newcommand{\R}{\mathbb{R}}
\newcommand{\N}{\mathbb{N}}
\newcommand{\C}{\mathbb{C}}
\newcommand{\lap}{\mathscr{L}}
\newcommand{\dd}{\mathrm{d}}
\newcommand{\smallmat}[1]{\left[ \begin{smallmatrix}#1 \end{smallmatrix} \right]}
\newcommand{\vv}[1]{\ensuremath{\underline{#1}}}
\newcommand{\mm}[1]{\ensuremath{\bm{#1}}}

%opening
\title[MATH 2860 (Elementary Differential Equations)]{Homework 12 for MATH 2860 (Elementary Differential Equations)}


\author[Emmanuel Atindama]{E. A. Atindama, PhD Mathematics}

\address[E. A. Atindama]{. Email addresses: \href{emmanuel.atindama@utoledo.edu}{emmanuel.atindama@utoledo.edu} 
% and 
% \href{mailto:a.student@qub.ac.uk}{a.student@qub.ac.uk}.
}
\thanks{Version 0.0.1. Last updated:~\today.}
\begin{document}

\maketitle

Due at 10:00 EST in class

\begin{enumerate}[1.]
  \item Convert the IVP
        \[
        y'' + 4y = \sin t, \quad y(0)=0, \ y'(0)=1
        \]
        into a first-order system. \textcolor{red}{do not solve!}
  
  \item Given the system of algebraic equations 
      \[
      \begin{aligned}
        x + 3y +z  &= 1 \\
        1 + z + 2y &= 0 \\
        y+z+x &= 1
      \end{aligned}
      \]
      \begin{enumerate}[(a)]
      \item Convert the system into a matrix equation \(A \vec{v} = \vec{b}\), where \(\vec{v} =\begin{bmatrix} x & y & z \end{bmatrix}\) is the vector of unknowns.
        % \solspace{1.5in}
      \item Without finding solutions, perform a calculation that can determine if this matrix equation has a \emph{unique} solution vector?
      % \solspace{2in}
      \item Use Gauss elimination to find the value or a general formula for the solution vector.
      % \solspace{2in}
      \end{enumerate}

  \item Solve the nonhomogeneous system
      \[
          \mathbf{x}' = \begin{bmatrix} 1 & -1 \\ 2 & 0 \end{bmatrix} \mathbf{x} + \begin{bmatrix} e^t \\ t \end{bmatrix}, \quad \mathbf{x}(0) = \begin{bmatrix} 0 \\ 1 \end{bmatrix},
      \]
      using the \textbf{variation of constants formula}.


  \item Given is the second order differential equation \[2 \ddot y (t) + c \dot y(t) + 8 y(t) = 0,\] where \(c\) is a parameter.
      \begin{itemize}
      \item[(i)] Write the ODE as the first-order system of differential
        equations. \footnote{Introduce variables \(x_{1} = y\) and \(x_{2} =
          \dot{y}\).}
      \item[(ii)] Compute the eigenvalues of the obtained system (they will depend on \(c\)).
      \item[(iii)] Find the range of values of \(c\) that makes the equilibrium at origin \emph{stable} and that results in oscillating trajectories (spiral phase portrait).
    \end{itemize}

\end{enumerate}

  \textcolor{red!20}{Show all your work.}
\end{document}
